Recent advances in large language models have enabled increasingly capable AI assistants, yet deploying these agents in real-time communication platforms like Discord presents unique challenges. Users expect assistants to be helpful and proactive, but also safe and accurate. Unproactive agents miss opportunities to assist, while overly aggressive agents can annoy users or perform unwanted actions. Furthermore, language models are prone to hallucinations and may lack up-to-date or domain-specific knowledge needed for effective assistance in specialized communities.

Existing approaches to AI agents in communication platforms typically fall into two categories: reactive bots that only respond to explicit commands, or proactive agents that initiate actions without sufficient user oversight. Reactive agents fail to anticipate user needs, while proactive agents without oversight can compromise safety and accuracy. Recent work on tool use in language models (Schick et al., 2023; Yao et al., 2023) shows promise for extending model capabilities, but often lacks the safety mechanisms needed for real-world deployment. Similarly, retrieval-augmented generation approaches improve factual accuracy but may not address the proactive assistance and safety concerns inherent in interactive environments.

We propose Discord OpenClaw, a framework that integrates three key components: (1) proactive behavior through periodic system monitoring and task suggestion, (2) structured approval workflows that require user consent before executing potentially impactful actions, and (3) external reasoning integration via tools like NotebookLM to ground responses in verified knowledge sources. This combination addresses the tension between helpfulness and safety, allowing agents to be proactive while maintaining user control and improving factual accuracy.

Our contributions are threefold: First, we introduce a comprehensive framework for safe, proactive AI agents in communication platforms that balances initiative with oversight. Second, we provide empirical evidence from controlled experiments demonstrating that this framework significantly improves user engagement, safety, trust, and answer quality. Third, we show through ablation studies that each component contributes uniquely to overall performance, with proactive behavior and approval workflows having the largest impact. To our knowledge, this is the first work to systematically study and evaluate this combination of techniques in the context of Discord-based AI agents.